\documentclass{article}
\usepackage{makecell}

\title{Ch. 5 Optional Problems}
\author{Principles of Microeconomics}
\date{\today}

\begin{document}

\maketitle

\begin{enumerate}

\item Fill in the missing amounts in the following table:

\begin{tabular}{| c | c | c | c |}
\hline
 & \makecell{\% CHANGE\\IN PRICE} & \makecell{\% CHANGE\\IN QUANTITY} & ELASTICITY \\
 \hline
 \makecell{Demand for Ben \& Jerry's\\Ice Cream} & +10\% & -12\% & \textbf{a.} \\
 \hline
 \makecell{Demand for beer at San\\Francisco 4934s football games} & -20\% & \textbf{b.} & 0.5 \\
 \hline
 \makecell{Demand for Broadway theater\\tickets in New York} & \textbf{c.} & -15\% & 1.0 \\
 \hline
 Supply of chickens & +10\% & \textbf{d.} & 1.2 \\
 \hline
 Supply of beef cattle & -15\% & -10\% & \textbf{e.} \\
 \hline
\end{tabular}

\item Use the table in the preceding problem to defend your answers to the following questions:

	\begin{enumerate}
	
	\item Would you recommend that Ben \& Jerry's move forward with a plan to raise prices if the company's only goal is to increase revenues?
	
	\item Would you recommend that beer stands cut prices to increase revenues at 49ers football games next year?
	
	\end{enumerate}

\item For each of the following scenarios, decide whether you agree or disagree and explain your answer.

	\begin{enumerate}
	
	\item If the elasticity of demand for cocaine is 0.2 and the Drug Enforcement Administration succeeds in reducing supply substantially, causing the street price of the drug to rise by 50\%, buyers will spend less on cocaine.
	
	\item Every year Christmas tree vendors bring tens of thousands of trees from the forests of New England to New York City and Boston. During the last two years, the market has been very competitive; as a result, price has fallen by 10 percent. If the price elasticity of demand was 1.3, vendors would lose revenues altogether as a result of the price decline.
	
	\item If the demand for a good has unit elasticity, it is always true that an increase in its price will lead to more revenues for sellers taken as a whole.
	
	\end{enumerate}

\item Taxicab fares in most cities are regulated. Several years ago taxi-cab drivers in Boston obtained permission to raise their fares 10 percent, and they anticipated that revenues would increase by about 10 percent as a result. They were disappointed, however. When the commissioner granted the 10 percent increase, revenues increased by only about 5 percent. What can you infer about the elasticity of demand for taxicab rides? What were taxicab drivers assuming about the elasticity of demand?

\item The cross price elasticity values for three sets of products are listed in the table below. What can you conclude about the relationships between each of these sets of products?

\begin{tabular}{| c | c | c | c |}
\hline
& \makecell{PRODUCTS\\A \& B} & \makecell{PRODUCTS\\C \& D} & \makecell{PRODUCTS\\E \& F} \\
\hline
\makecell{Cross-price\\elasticity} & -8.7 & 5.5 & 0.0 \\
\hline
\end{tabular}

\end{enumerate}

\end{document}